% The spoken body of the homily for the Twentieth Sunday after Pentecost.
% Continuous preaching, ready to read aloud: no heading, no direction to the
% preacher, no citation inside the speech. The loci, the relation to the
% reviewed readings, the audience, the word count and the pace estimate stand
% in the terminal note, which no one reads out. The blank lines are the
% paragraph breaks a speaker recovers his place by, and the vertical spaces
% divide the speech into its movements: the man on the road with a sentence;
% what he asked, the rebuke, and the Fathers' two weights for it; the turn,
% the plan refused and the prayer granted, with Gregory's command; the walk
% and the night, with Aquinas and the Gradual's due season; the prayers not
% answered as asked, with Chrysostom, the Epistle and the Communion's unsung
% line; the father's place at this altar, with the Secret; the
% Introit, the Collect and the Postcommunion; the week's response, with the
% Offertory's exiles; and the road's end. The columns are set ragged at the
% foot, so that the breaks keep one height instead of absorbing the columns'
% balancing.
\raggedcolumns

``Go thy way. Thy son liveth.''

Those words were all a man was given to carry home. He was a royal official
from Capharnaum, and his son was dying of a fever. He went to find Jesus at
Cana, and he went with a plan: Jesus would come down with him, stand by the
bed, and heal the boy. He came home without Jesus, with nothing in his hands
but a sentence.

That is today's Gospel, and it asks a question many of us have lived. What do
you do when you have asked God to come, and he sends you a word instead?

\medskip

Listen to what the man asked. ``Lord, come down before that my son die.'' The
first answer he receives is hard. ``Unless you see signs and wonders, you
believe not.''

The Fathers did not agree how hard to take that. St Augustine took it at full
weight: the Lord, he says, ``shows us a man lukewarm, or cold in faith, or of
no faith at all.'' St Gregory the Great is gentler. No one, he says, asks to be
saved by someone he does not believe can save; so the man believed. But he
believed and doubted. He asked the Lord to come down. He was looking for the
Lord's bodily presence, when the Lord is absent nowhere.

St John Chrysostom hears mercy in the rebuke: ``Here He healeth the father,
sick in mind, no less than the son.'' There were two patients in this story.
One lay in bed at Capharnaum. The other stood in front of Jesus, asking.

\medskip

So the man asks again: ``Lord, come down.'' And Jesus does not come down. He
says, ``Go thy way. Thy son liveth.''

Christ refuses the man's plan, and he grants the man's prayer. The father asked
for two things, a presence and a life. He is given the life at once. And he is
given the presence too, though not in the shape he imagined. He is given it as
a word. St Gregory says that the one who made all things by his will gave the
boy health by his command alone. That word was not less than Christ's coming.
It was the power that made the world, and for that power there is no distance
between Cana and Capharnaum.

This is what the Gospel changes in us. We tend to think that where we cannot
see God, he is absent; that if he has not come the way we asked, he has not
come. The father thought so too. But what looked like Christ staying away was
Christ at work, by a word.

\medskip

Then comes one of the bravest sentences in Scripture. ``The man believed the
word which Jesus said to him and went his way.''

He went. He had no proof. For all he knew, the fever was still burning. Faith
here is not a feeling. It is walking on a word. St Thomas Aquinas says that on
God's road not to go forward is to go back. This man's faith grew as he walked.

His servants met him on the road and told him, ``Yesterday at the seventh hour,
the fever left him.'' Yesterday. The word had done its work before he knew it.
He spent a night between the word and the news.

The Gradual speaks of such a night. ``The eyes of all hope in thee, O Lord: and
thou givest them meat in due season.'' St Augustine explains ``due season''
from the sickroom: whoever tends a sick man gives him food ``when he ought to
receive.'' And to anyone who has prayed and not yet seen the answer he says,
``let his eyes wait for the food, which He giveth in due season.... Thou
delayest, not deniest.''

\medskip

But we must be honest, because this Gospel is honest. Some here have prayed for
a son or a daughter, a husband or a wife, as this man prayed, and the answer
was not his. The Gospel does not promise that every prayer will be answered in
the shape we name. St John Chrysostom, preaching on this passage, saw people
give thanks when a child recovered. But they ought, he said, ``even if they
obtain it not, to persist just the same in giving thanks, in glorifying God.''
That is also the word of today's Epistle: ``Giving thanks always for all
things.''

How can we give thanks when we have not been given what we asked? Not because
loss is good. Because of who has spoken. The Communion prays, ``Be thou mindful
of thy word to thy servant, in which thou hast given me hope.'' God does not
need reminding, Augustine says; to ask him to remember is to say, ``fulfil Thy
promise to Thy servant.'' And the psalm goes on, in words the Communion leaves
out: ``because thy word hath enlivened me.'' Thy word hath given me life; thy
son liveth. The deepest promise of God's word is not that this fever will break
tomorrow. It is life, the life God gives that does not end. That promise stands
over every prayer we make, the prayers answered as we asked and the prayers
that were not.

\medskip

We are where the father stood. Many of us came here carrying someone, and have
said in one way or another: Lord, come down before it is too late. And the
Mass does what Christ did at Cana. It gives us a word, the Gospel we have
heard. Then, over the gifts, the priest will ask, ``May these mysteries, O
Lord, we beseech thee, procure us a heavenly remedy, and cleanse away the vices
of our hearts.'' The Latin word for that remedy is medicine. The one who heals
in these mysteries is the one who healed by a word at Capharnaum.

The Church believes even more than that father dared to ask. On this altar
Christ does come down: his Body and Blood are truly here. But he comes as he
came to Capharnaum, by a word, hidden from our eyes. What we see is bread and
wine. What we are asked is what the father was asked: to believe the word, and
to go our way.

\medskip

This Mass began with words from the prayer of Azarias in the furnace at
Babylon. ``Whatever thou hast done to us, O Lord, thou hast done by a just
judgment: for we have sinned and disobeyed thy commandments: but glorify thy
name, and deal with us according to thy great mercy.'' St Augustine says that
the saints pray in that order, God's justice first and his help after: ``First
they praised Him scourging, and so they felt Him feeding.'' The father had to
learn that order: not first my plan and then God, but first God, and then my
need, placed in his hands.

The Collect asked that we may ``serve thee with secure minds.'' ``Secure'' has
its old meaning here: free from care, at rest. It is the mind of that man on
the road in the dark. He was not sure what he would find. He was sure who had
spoken.

And the last prayer of this Mass will answer the first. We confessed that we
``disobeyed thy commandments.'' At the end we shall ask, ``grant, we beseech
thee, we may always obey thy commandments.'' Between the two stands the word we
receive. The word received becomes the word obeyed. ``Go thy way'' was a
command, and the father obeyed it.

\medskip

So here is something for this week. When you have prayed for someone, go your
way, as the father did. Do not stand at Cana, refusing to move until God acts
as you planned. Take the Communion verse with you, and say it each morning:
``Be thou mindful of thy word to thy servant, in which thou hast given me
hope.'' Then do the next thing God asks of you: the work in front of you, the
person in your house, the apology you owe. And give thanks before you see the
answer.

The Offertory will sing of exiles weeping by the rivers of Babylon, remembering
Sion. Preaching on that psalm, St Augustine told his people, ``whither your
hope goeth before, let your life follow.'' That is the father's road. His hope
went home ahead of him, carried by a word, and his feet followed.

\medskip

At the end of that road, the evangelist says, ``himself believed, and his whole
house.'' The faith he brought home was larger than the faith he set out with.
It filled a household.

We are on that road now, between the word and the sight. Christ has spoken, and
his word gives life. And the Gospel tells us what to do on the way: ``The man
believed the word which Jesus said to him and went his way.''
